%! Author = chtompki
%! Date = 1/14/26
% Example LaTeX document for GP111 - note % sign indicates a comment
\documentclass[12pt]{amsart}
% Default margins are too wide all the way around. I reset them here
\setlength{\hoffset}{-.75in}
\setlength{\textwidth}{6.5in}
\setlength{\voffset}{-.5in}
\setlength{\textheight}{9.0in}
\usepackage{amsmath}
\usepackage{amssymb}
\usepackage{pgfplots}
\usepackage{tikz}
\usepackage{amsthm}
\usepackage{enumitem}
\usepackage{setspace}

\newenvironment{solution}
{\begin{proof}[Solution]}
{\end{proof}}

\title{Advanced Linear Algebra\\Homework 6}
\author{Rob Tompkins}
\date{\today}

\begin{document}
    \maketitle
    \date{\today}
    \begin{onehalfspace}

        % P.4.9
        \noindent\textbf{P.4.9.} If $A\in\mathbf{M}_n$ is to be represented as $A=XY^\top$ for some
        $X,Y\in\mathbf{M}_{n\times r}$, explain why $r$ cannot be smaller than $\text{rank}(A)$.
        \begin{solution}
        \end{solution}

        % P.4.10
        \noindent\textbf{P.4.10.} For $X\in\mathbf{M}_{m\times n}$, $\nu(X)$ denote the nullity of $X$.
        \begin{enumerate}[label=(\alph*)]
            \item Show that $\nu(X^\top) = \nu(X) + m - n$.
            \item Let $A\in\mathbf{M}_{m\times k}$ and $B\in\mathbf{M}_{k\times n}$.
                  Show that the rank inequality given by
                \begin{displaymath}
                    \text{rank}(A) + \text{rank}(B) - k \le \text{rank}(AB)
                \end{displaymath}
                  is equivalent to
                \begin{displaymath}
                    \nu(AB) \le \nu(A) + \nu(B).
                \end{displaymath}
            \item If $A,B\in\mathbf{M}_{n}$, show that the rank inequality given by
            \begin{displaymath}
                \text{rank}(AB) \le \text{min}\{\text{rank}(A), \text{rank}(B)\}.
            \end{displaymath}
            is equivalent to
        \end{enumerate}
        \begin{solution}
        \end{solution}

        % P.4.13
        \noindent\textbf{P.4.13.}
        \begin{solution}
        \end{solution}

        % P.4.35
        \noindent\textbf{P.4.35.}
        \begin{solution}
        \end{solution}

        % P.4.38
        \noindent\textbf{P.4.38.}
        \begin{solution}
        \end{solution}

    \end{onehalfspace}
\end{document}
